\section{Краткая характеристика пользовательского интерфейса}
\subsection{Задание параметров}
Начальное окно содержит группы «Компания и время», «Сценарий», «Периодические семинары», «Напоминания», «Конфликты и эксперимент». Задаются размер компании, период, шаг, рабочая неделя, интенсивности, участники, сроки, важность, роль и seed. Не относящиеся к выбранной политике поля поиска отключаются.

Доступны проверка параметров, возврат к значениям по умолчанию, сохранение и загрузка настроек, открытие состояния и кнопка «Провести серию экспериментов…». Кнопка «Открыть моделирование» создаёт независимый прогон с копией настроек. Изменение исходной формы не меняет параметры уже открытого движка.

\subsection{Окно моделирования}
В верхней части расположены запуск, пауза, следующий шаг, сброс, скорость и «+ Событие». Сводка показывает модельное время и количество активных и архивных событий. Значения скорости 0.5$\times$, 1$\times$, 2$\times$ и 5$\times$ регулируют таймер и анимацию.

\begin{longtable}{L{0.25\linewidth}L{0.67\linewidth}}
\toprule Вкладка & Назначение \\\midrule\endhead
Журнал & Уже совершённые действия с модельным временем. Двойной щелчок открывает соответствующее мероприятие. \\
Календарь и архив & Просмотр карточек по датам, статусам, сотрудникам, отделам и часам внутри каждого дня. Завершённые мероприятия доступны в архиве. \\
Офис & План помещений, рабочие места и зоны встреч. Движущиеся объекты показывают переходы сотрудников между местами работы, мероприятиями и командировками. \\
Участники & Публичный ID, имя, отдел и текущая занятость; руководитель дополнительно обозначен. \\
Отделы & Название, число сотрудников и руководитель. Ручных кадровых кнопок нет. \\
Типы событий & Каталог пяти типов мероприятий и операции переноса с весами и длительностями. \\
Заявки, почта и отчёт & Рассмотрение конфликтов, имитируемые сообщения, полный список событий, сводка, сохранение и экспорт. \\
\bottomrule
\end{longtable}

\subsection{Добавление, редактирование и конфликты}
Форма мероприятия содержит название, тип, начало и конец, помещение, участников, важность и организатора. Участники задаются публичными ID через запятую. После сохранения ядро проверяет данные, права, занятость и пересечения. Некорректный ввод не создаёт подтверждённое мероприятие.

Открытие формы ставит прогон на паузу. Запланированное событие редактируется только при наличии полномочий. Начатые, завершённые и недоступные по роли события показываются без возможности сохранения. Отдельной кнопки отмены мероприятия нет.

При \code{Suggest} новые ожидающие заявки останавливают автоматическое продвижение. На вкладке «Заявки» показываются причина и ID конфликтующих событий. Кнопка «Выбрать вариант…» открывает предложенные время и помещения; выбор проверяется повторно. «Редактировать заявку…» снимает ожидание старой попытки и открывает её поля для новой. Для изменения уже существующего конфликтующего мероприятия пользователь может открыть его карточку в календаре в пределах своей роли.

\subsection{Почта, отчёт и исследование}
Имитируемые сообщения объединяются для просмотра по рассылке; полный текст и адресаты открываются отдельно. В статистике доставка считается каждому адресату. Экспорт JSON/CSV доступен как до завершения, так и после него; промежуточный отчёт содержит текущее модельное время.

Окно серии экспериментов показывает число завершённых запусков и агрегаты стратегий. Остановленная серия отмечается незавершённой. В отличие от обычного окна, выбор в пакетном \code{Suggest} выполняет программное правило, поэтому серия не ждёт участия человека.
